\section{The box ensemble: root location and coefficient scale}\label{sec:box}

Fix $d,n\ge1$ and let
\begin{equation}\label{eq:box-model}
 p_i(z)=\sum_{\alpha\in\{0,\ldots,n\}^d}c_{i,\alpha}z^\alpha,
 \qquad c_{i,\alpha}\sim\CN(0,1)\quad\hbox{independently}.
\end{equation}
The monomials are orthonormal in the coefficient norm. For generic coefficients,
the system has
\begin{equation}\label{eq:box-count}
 N=d!n^d
\end{equation}
simple affine roots, all in $(\C^*)^d$. This is the mixed-volume count for
$d$ identical boxes $[0,n]^d$ \cite{bernstein1975roots,cox2005using}.
The exceptional coefficient sets are proper algebraic sets and hence have
Gaussian measure zero. Equivalently, the multihomogeneous compactification
on $(\mathbb P^1)^d$ has intersection number $d!n^d$, with no boundary
intersection for generic sections. This statement concerns that
compactification; ordinary homogenization to total degree $dn$ can introduce
common zeros at projective infinity.

Put
\[
 k_n(r)=\sum_{j=0}^n r^j,\qquad
 h_n(r)=\left(r\frac{k_n'(r)}{k_n(r)}\right)'.
\]
The covariance kernel and its metric are
\begin{equation}\label{eq:box-kernel}
 K(z,w)=\prod_{\ell=1}^d k_n(z_\ell\bar w_\ell),\qquad
 H(z)=\diag(h_n(|z_1|^2),\ldots,h_n(|z_d|^2)).
\end{equation}
In particular $H>0$. A convenient interpretation is obtained by giving
$J_r\in\{0,\ldots,n\}$ probabilities
$\Prob(J_r=j)=r^j/k_n(r)$. Differentiation gives
\begin{equation}\label{eq:box-variance}
 r h_n(r)=\Var(J_r),\qquad r>0.
\end{equation}
The continuous endpoint values are $h_n(0)=1$ and
$h_n(1)=n(n+2)/12$.

\subsection{An exact selected-root law}

\begin{proposition}[Independent radial coordinates]\label{prop:box-root-law}
Under $\Pstar$, the coordinates of $Z$ are independent, their angles are
independent and uniform, and $R_\ell=|Z_\ell|^2$ has distribution
\begin{equation}\label{eq:box-cdf}
 F_n(r)=\Pstar(R_\ell\le r)=\frac{r k_n'(r)}{n k_n(r)},\qquad
 F_n'(r)=\frac{h_n(r)}n\quad(r\ge0).
\end{equation}
Moreover, $R_\ell$ and $R_\ell^{-1}$ have the same law.
\end{proposition}
\begin{proof}
Theorem~\ref{thm:local-universal} and \eqref{eq:box-count} give joint density
$\prod_\ell h_n(|z_\ell|^2)/(n\pi)$ with respect to complex Lebesgue measure.
In polar coordinates $\dd V(z_\ell)=\tfrac12\dd r_\ell\dd\theta_\ell$,
which proves the density statement. Integration gives \eqref{eq:box-cdf},
with limits zero and one since $r k_n'/k_n$ tends to zero and $n$ at the
two endpoints. The identity $k_n(r)=r^n k_n(r^{-1})$ implies
$F_n(r)=1-F_n(r^{-1})$ for $r>0$, proving inversion symmetry.
\end{proof}

The independence in this proposition is between coordinates of one selected
root. Different roots of the same polynomial system remain dependent.

\begin{lemma}[Depth and the outer-root moment]\label{lem:box-depth}
For any one coordinate, put $s=n\log R$. Then
\begin{equation}\label{eq:box-depth}
 \Pstar(|s|>t)\le\min\{1,2/t\}\quad(t>0),\qquad
 \Estar\log^+R=\frac{\log(n+1)}n.
\end{equation}
Consequently $\sup_n\Estar\log(1+|s|)<\infty$.
\end{lemma}
\begin{proof}
For $0<r<1$, the mean of the truncated geometric law in
\eqref{eq:box-variance} is at most the untruncated mean $r/(1-r)$.
Thus
\[
 F_n(e^{-t/n})\le\frac1{n(e^{t/n}-1)}\le\frac1t.
\]
Inversion symmetry gives the two-sided tail. Also
\begin{align*}
 \Estar\log^+R
 &=\int_1^\infty\frac{1-F_n(r)}r\dd r\\
 &=\frac1n\left[\log\frac{r^n}{k_n(r)}\right]_{r=1}^{r=\infty}
 =\frac{\log(n+1)}n.
\end{align*}
Finally integrate the depth tail against $(1+t)^{-1}\dd t$.
\end{proof}

\subsection{Uniform metric estimates}

The next elementary estimate is used both for coefficient conditioning and
for the companion overlap. It replaces a boundary-layer asymptotic by a bound
valid at every radius and degree.

\begin{lemma}[The one-variable kernel metric]\label{lem:box-metric-bound}
There are absolute constants $0<c<C<\infty$ such that, for $n\ge1$ and
$0\le r\le1$,
\begin{equation}\label{eq:box-metric-bound}
 \frac{c}{(1-r+1/n)^2}\le h_n(r)
 \le\frac{C}{(1-r+1/n)^2}.
\end{equation}
For $r>0$,
\begin{equation}\label{eq:box-metric-reversal}
 h_n(r)=r^{-2}h_n(r^{-1}).
\end{equation}
\end{lemma}
\begin{proof}
For $0\le r\le1/4$, the pair of indices $0,1$ in the variance formula gives
$h_n(r)\ge k_n(r)^{-2}\ge(1-r)^2\ge9/16$. Direct differentiation of
$k_n(r)=(1-r^{n+1})/(1-r)$ gives, for $0<r<1$,
\begin{equation}\label{eq:box-explicit-h}
 h_n(r)=\frac1{(1-r)^2}
       -\frac{(n+1)^2r^n}{(1-r^{n+1})^2}
 \le\frac1{(1-r)^2}.
\end{equation}
Thus $h_n$ is bounded above and below by absolute positive constants on
$[0,1/4]$, which proves the assertion there.

For $1/4\le r<1$, set
\[
 m=\min\{n,\lfloor(1-r)^{-1}\rfloor\}.
\]
Then $k_n(r)\le m+1$. Also $\log r\ge-(1-r)/r\ge-4(1-r)$, so
$r^j\ge e^{-4}$ for $0\le j\le m$. Each of these indices has probability
at least $e^{-4}/(m+1)$. Using an independent copy $J_r'$,
\[
 \Var(J_r)=\frac12\E(J_r-J_r')^2
 \ge e^{-8}\Var(\operatorname{Unif}\{0,\ldots,m\})
 =e^{-8}\frac{m(m+2)}{12}.
\]
Since $m\ge\tfrac12\min\{n,(1-r)^{-1}\}$, division by $r\le1$
proves the lower bound in \eqref{eq:box-metric-bound}.
For the upper bound, $\Var(J_r)\le n^2/4$ and $r\ge1/4$ imply
$h_n(r)\le n^2$. Combining this with \eqref{eq:box-explicit-h} gives
\[
 h_n(r)\le\min\{n^2,(1-r)^{-2}\}
 \le\frac4{(1-r+1/n)^2}.
\]
At $r=1$, use $h_n(1)=n(n+2)/12$. Finally, differentiating
$F_n(r)=1-F_n(r^{-1})$ proves \eqref{eq:box-metric-reversal}.
\end{proof}

\begin{corollary}[Logarithmic metric scale]\label{cor:box-metric-log}
For $R$ with distribution \eqref{eq:box-cdf},
\begin{equation}\label{eq:box-metric-log}
 \Estar\left|\log\frac{h_n(R)}{n^2}\right|\le C
\end{equation}
with an absolute constant. In particular $\Estar\log h_n(R)=2\log n+O(1)$.
\end{corollary}
\begin{proof}
For $0<r\le1$, write $s=n\log r\le0$. Since
$n(1-r)\le|s|$, Lemma~\ref{lem:box-metric-bound} gives
\[
 \frac{c}{(1+|s|)^2}\le\frac{h_n(r)}{n^2}\le C.
\]
For $r>1$, apply reversal. At every $r>0$ this yields
\[
 \left|\log\frac{h_n(r)}{n^2}\right|
 \le C+2\log(1+|n\log r|)+2\log^+r.
\]
Take expectations and apply Lemma~\ref{lem:box-depth}, noting that
$\log(n+1)/n\le\log2$.
\end{proof}

\subsection{Coefficient conditioning and representation overlap}

\begin{proposition}[The coefficient scale]\label{prop:box-coefficient-scale}
For fixed $d$ in the box model,
\begin{equation}\label{eq:box-coefficient-scale}
 \Estar\log\kappa_{\mathrm{coeff}}(P,Z)=-\log n+O_d(1).
\end{equation}
The same formula holds for the matched scalar coordinate condition
$\kappa_{\mathrm{coeff}}^{x_1}$ in \eqref{eq:setup-scalar-condition}.
All logarithms here are absolutely integrable for each $n,d$.
\end{proposition}
\begin{proof}
Put $H=n^2D$, with diagonal entries $h_n(R_\ell)/n^2$. From
$J=\sqrt S A H^{1/2}$ we obtain
\[
 \kappa_{\mathrm{coeff}}=n^{-1}\|D^{-1/2}A^{-1}\|_2.
\]
The logarithm of $\|D^{-1/2}A^{-1}\|_2$ differs from that of
$\|A^{-1}\|_2$ by at most
$\tfrac12\max_\ell|\log D_{\ell\ell}|$.
Theorem~\ref{thm:local-tail} and Corollary~\ref{cor:box-metric-log}
give \eqref{eq:box-coefficient-scale} and absolute integrability.
For the scalar output there is the more precise factorization
\begin{equation}\label{eq:box-coordinate-factorization}
 \kappa_{\mathrm{coeff}}^{x_1}
 =h_n(R_1)^{-1/2}\|e_1^T A^{-1}\|_2.
\end{equation}
The root location and $A$ are independent. The row norm has an integrable
logarithm because it is at most $\|A^{-1}\|_2$ and at least
$\|A\|_2^{-1}$. Apply Corollary~\ref{cor:box-metric-log} again.
\end{proof}

Suppose now that a fixed monomial set $\B$ is a quotient basis on a generic
coefficient chart. Use multiplication by $x_1$, whose root values are distinct
almost surely in the box model. One way to check the last assertion is to
consider pairs $z\ne w$ with $z_1=w_1$: evaluation at these points gives two
independent linear conditions on each polynomial, since constants and all
coordinates belong to the coefficient space. The pair incidence set has
dimension at most one less than coefficient space. Its projection therefore
has Gaussian measure zero.

Define the extended nonnegative expectation
$\Theta_{n,d}^{\B}=\Estar\log\Lambda_{\B}$, where the overlap uses the
Euclidean norm on the chosen quotient coordinates. Equations
\eqref{eq:setup-overlap} and \eqref{eq:box-coefficient-scale} imply
\begin{equation}\label{eq:box-decomposition}
 \Estar\log\frac{\kappa_{\mathrm{eig}}(M_{x_1},Z_1)}
                     {\kappa_{\mathrm{coeff}}^{x_1}(P,Z)}
 =\log n+\Theta_{n,d}^{\B}+O_d(1).
\end{equation}
This identity also holds with the full-root condition in the denominator,
but \eqref{eq:box-decomposition} compares the same scalar output.
If $\Theta_{n,d}^{\B}=+\infty$, interpret the equality in the extended sense;
the other logarithmic terms are integrable. No bound on multivariate overlap
is asserted here. Section~\ref{sec:companion} controls it for the univariate
companion basis.

For completeness, the evaluation factor itself satisfies
\begin{equation}\label{eq:box-evaluation-scale}
 \Estar\log k_n(R)=2\log n+O(1),\qquad
 \Estar\log\sqrt{S(Z)}=d\log n+O_d(1).
\end{equation}
Indeed reversal gives the exact identity
\[
 \log k_n(R)=\log k_n(e^{-|s|/n})+s_+.
\]
Here $\Estar s_+=\log(n+1)$, while
\[
 \frac{n}{1+|s|}\le k_n(e^{-|s|/n})\le n+1.
\]
The lower bound follows by comparing the sum with
$\int_0^n e^{-|s|x/n}\dd x$ and using $1-e^{-u}\ge u/(1+u)$.
Lemma~\ref{lem:box-depth} proves \eqref{eq:box-evaluation-scale}.
Replacing coefficient conditioning by residual conditioning would therefore
change the leading term in \eqref{eq:box-decomposition} from $\log n$ to
$(d+1)\log n$. The outer roots contribute to this difference, even though
most roots lie close to the unit polytorus.
